% GENERATED FILE. Edit canonical lessons, then run npm run generate:books.
% canonical-source-hash: fnv1a64:933452a66eb18eee
% canonical-lessons: ZH-C430-guanguang, ZH-C430-guji, ZH-C430-jingqu, ZH-C430-zipai, ZH-C430-wangji

\chapter{Seeing the Sights}
\label{ch:zh-seeing-the-sights}
\hlchaptermodality{430}

\begin{chapteropening}
\textbf{By the end of this chapter:} \emph{I can talk about going sightseeing, historic sites, a scenic area, a selfie and the busy season.}

Complete the last lesson of chapter 430: I can talk about going sightseeing, historic sites, a scenic area, a selfie and the busy season.
\end{chapteropening}

\section[guānguāng]{guānguāng — to go sightseeing}

\label{lesson:ZH-C430-guanguang}



\begin{quote}
Before the new one: say the Mandarin for to get about, then the Mandarin for travel costs.
\end{quote}

\subsection*{You'll want to know: guānguāng}
\emph{guānguāng} — \textquotedblleft{}to go sightseeing\textquotedblright{}.

\begin{quote}
\emph{wǒmen zhù zài zhè jiā jiǔdiàn měitiān qù guānguāng} — We're staying at this hotel and going sightseeing every day.
\end{quote}

Its characters come later; for now it is a word for the ear.

\begin{cousinweb}
\emph{guānguāng} is built from \emph{guān} 'to view' and \emph{guāng} 'light', both inherited from Middle Chinese, \emph{kwan} and \emph{kwang} (Guangyun, in Baxter's transcription), which glosses the first 'to view' and the second 'bright'. How the pair came to mean touring the sights is not settled.
\end{cousinweb}

\begin{grammarlens}[title={What you've built}]
One more word for getting around.
\end{grammarlens}

\subsection*{Your turn}
\begin{itemize}
\raggedright
  \item \emph{Say it:} \emph{guānguāng}
  \item \emph{Say it:} \emph{guānguāng}, once more
  \item \emph{Say it:} \emph{wǒmen zhù zài zhè jiā jiǔdiàn měitiān qù guānguāng} — We're staying at this hotel and going sightseeing every day.
  \item \emph{Recall:} say the Mandarin for to get about, then the Mandarin for travel costs, then say \emph{guānguāng} again
\end{itemize}

\begin{tcolorbox}[breakable,colback=teal!4,colframe=teal!35!black,title={Before you move on}]
What does guānguāng mean? (\textquotedblleft{}To go sightseeing\textquotedblright{}.) Say it once more.
\end{tcolorbox}
\section[gǔjì]{gǔjì — historic site}

\label{lesson:ZH-C430-guji}



\begin{quote}
Before the new one: say the Mandarin for travel costs, then the Mandarin for to go sightseeing.
\end{quote}

\subsection*{You'll want to know: gǔjì}
\emph{gǔjì} — \textquotedblleft{}historic site\textquotedblright{}.

\begin{quote}
\emph{wǒmen zuò gōngjiāo qù kàn gǔjì} — We're taking the bus to visit the historic sites.
\end{quote}

Its characters come later; for now it is a word for the ear.

\begin{cousinweb}
\emph{gǔjì} is built from \emph{gǔ} 'ancient' and \emph{jì} 'a trace', both inherited from Middle Chinese, \emph{kuX} and \emph{tsjek} (Guangyun, in Baxter's transcription), which glosses the first 'former' and the second 'a footprint'.
\end{cousinweb}

\begin{grammarlens}[title={What you've built}]
One more word for getting around.
\end{grammarlens}

\subsection*{Your turn}
\begin{itemize}
\raggedright
  \item \emph{Say it:} \emph{gǔjì}
  \item \emph{Say it:} \emph{gǔjì}, once more
  \item \emph{Say it:} \emph{wǒmen zuò gōngjiāo qù kàn gǔjì} — We're taking the bus to visit the historic sites.
  \item \emph{Recall:} say the Mandarin for travel costs, then the Mandarin for to go sightseeing, then say \emph{gǔjì} again
\end{itemize}

\begin{tcolorbox}[breakable,colback=teal!4,colframe=teal!35!black,title={Before you move on}]
What does gǔjì mean? (\textquotedblleft{}Historic site\textquotedblright{}.) Say it once more.
\end{tcolorbox}
\section[jǐngqū]{jǐngqū — scenic area}

\label{lesson:ZH-C430-jingqu}



\begin{quote}
Before the new one: say the Mandarin for to go sightseeing, then the Mandarin for a historic site.
\end{quote}

\subsection*{You'll want to know: jǐngqū}
\emph{jǐngqū} — \textquotedblleft{}scenic area\textquotedblright{}.

\begin{quote}
\emph{zhège jǐngqū zǎoshang bā diǎn kāimén} — This scenic area opens at eight in the morning.
\end{quote}

Its characters come later; for now it is a word for the ear.

\begin{cousinweb}
\emph{jǐngqū} is built from \emph{jǐng} 'a view' and \emph{qū} 'an area', both inherited from Middle Chinese, \emph{kjaengX} and \emph{khju} (Guangyun, in Baxter's transcription), which glosses the first 'bright' and 'great' and gives the second as the name of a marsh; the senses 'a view' and 'an area' are not given there.
\end{cousinweb}

\begin{grammarlens}[title={What you've built}]
One more word for getting around.
\end{grammarlens}

\subsection*{Your turn}
\begin{itemize}
\raggedright
  \item \emph{Say it:} \emph{jǐngqū}
  \item \emph{Say it:} \emph{jǐngqū}, once more
  \item \emph{Say it:} \emph{zhège jǐngqū zǎoshang bā diǎn kāimén} — This scenic area opens at eight in the morning.
  \item \emph{Recall:} say the Mandarin for to go sightseeing, then the Mandarin for a historic site, then say \emph{jǐngqū} again
\end{itemize}

\begin{tcolorbox}[breakable,colback=teal!4,colframe=teal!35!black,title={Before you move on}]
What does jǐngqū mean? (\textquotedblleft{}Scenic area\textquotedblright{}.) Say it once more.
\end{tcolorbox}
\section[zìpāi]{zìpāi — to take a selfie}

\label{lesson:ZH-C430-zipai}



\begin{quote}
Before the new one: say the Mandarin for a historic site, then the Mandarin for a scenic area.
\end{quote}

\subsection*{You'll want to know: zìpāi}
\emph{zìpāi} — \textquotedblleft{}to take a selfie\textquotedblright{}.

\begin{quote}
\emph{tā zài jiǔdiàn de dàtáng zìpāi} — She's taking a selfie in the hotel lobby.
\end{quote}

Its characters come later; for now it is a word for the ear.

\begin{cousinweb}
\emph{zìpāi} is built from \emph{zì} 'oneself' and \emph{pāi} 'to snap a picture', both inherited from Middle Chinese, \emph{dzijH} and \emph{phaek} (Guangyun, in Baxter's transcription), which glosses the first 'from' and the second 'to strike'; the sense 'to take a photo' is modern.
\end{cousinweb}

\begin{grammarlens}[title={What you've built}]
One more word for getting around.
\end{grammarlens}

\subsection*{Your turn}
\begin{itemize}
\raggedright
  \item \emph{Say it:} \emph{zìpāi}
  \item \emph{Say it:} \emph{zìpāi}, once more
  \item \emph{Say it:} \emph{tā zài jiǔdiàn de dàtáng zìpāi} — She's taking a selfie in the hotel lobby.
  \item \emph{Recall:} say the Mandarin for a historic site, then the Mandarin for a scenic area, then say \emph{zìpāi} again
\end{itemize}

\begin{tcolorbox}[breakable,colback=teal!4,colframe=teal!35!black,title={Before you move on}]
What does zìpāi mean? (\textquotedblleft{}To take a selfie\textquotedblright{}.) Say it once more.
\end{tcolorbox}
\section[wàngjì]{wàngjì — peak season; busy season}

\label{lesson:ZH-C430-wangji}



\begin{quote}
Before the new one: say the Mandarin for a scenic area, then the Mandarin for to take a selfie.
\end{quote}

\subsection*{You'll want to know: wàngjì}
\emph{wàngjì} — \textquotedblleft{}peak season; busy season\textquotedblright{}.

\begin{quote}
\emph{wàngjì de shíhou jiǔdiàn dōu kèmǎn le} — In the peak season the hotels are all fully booked.
\end{quote}

Its characters come later; for now it is a word for the ear.

\begin{cousinweb}
\emph{wàngjì} is built from \emph{wàng} 'thriving' and \emph{jì} 'a season', both inherited from Middle Chinese, \emph{hjwangH} and \emph{kjwijH} (Guangyun, in Baxter's transcription), which glosses the first as a fine light and the second as the youngest of brothers; the senses 'thriving' and 'a season' are not given there.
\end{cousinweb}

\begin{grammarlens}[title={What you've built}]
One more word for getting around.
\end{grammarlens}

\subsection*{Your turn}
\begin{itemize}
\raggedright
  \item \emph{Say it:} \emph{wàngjì}
  \item \emph{Say it:} \emph{wàngjì}, once more
  \item \emph{Say it:} \emph{wàngjì de shíhou jiǔdiàn dōu kèmǎn le} — In the peak season the hotels are all fully booked.
  \item \emph{Recall:} say the Mandarin for a scenic area, then the Mandarin for to take a selfie, then say \emph{wàngjì} again
\end{itemize}

\begin{tcolorbox}[breakable,colback=teal!4,colframe=teal!35!black,title={Before you move on}]
What does wàngjì mean? (\textquotedblleft{}Peak season; busy season\textquotedblright{}.) Say it once more.
\end{tcolorbox}
